
\section{Formulación del problema}
¿Cómo puede desarrollarse una metodología de ciencia de datos, basada únicamente en información transaccional parcial, que permita estimar la \emph{participación en la billetera} (SoW) y el \emph{tamaño de la billetera} (SioW) de los clientes empresariales de una empresa distribuidora de biblias y literatura cristiana, caracterizar su comportamiento de compra y apoyar su análisis mediante segmentación exploratoria, en ausencia de datos de encuestas y de transacciones con competidores?

Las preguntas de sistematización que apoyan la pregunta de formulación:

\begin{enumerate}
    \item ¿Qué variables transaccionales y de perfilamiento del cliente, identificadas mediante modelos estadísticos y técnicas de ciencia de datos aplicadas a información transaccional parcial, resultan relevantes para explicar la \emph{participación en la billetera} (SoW) y el \emph{tamaño de la billetera} (SioW) de los clientes empresariales de la empresa focal, en ausencia de información externa sobre competidores?

    \item ¿Cómo afecta la especificación del conjunto de covariables transaccionales —utilizando variables originales frente a una representación reducida mediante Análisis de Componentes Principales (PCA)— la estabilidad, parsimonia e interpretabilidad de las estimaciones de SoW y SioW obtenidas a partir del modelo \emph{pogit}?

    \item ¿En qué medida los segmentos de clientes empresariales identificados mediante técnicas de aprendizaje no supervisado (k-means y k-medoids) presentan correspondencia con los patrones emergentes en el espacio conjunto SoW--SioW estimado por el modelo \emph{pogit}, y cómo contribuye esta correspondencia a la interpretación económica de los resultados y a la toma de decisiones estratégicas de marketing?
\end{enumerate}
